\begin{theorem}[Existence and uniqueness]\label{thm:main}
Let $n=6$. Then $\det G_n(c)=\left(1 - c\right)^{5} \left(5 c + 1\right)$ and $G_n(c)$ is positive definite if and only if $- \frac{1}{5}<c<1$.
For every such $c$ there is a cubic rhombus $R_n(c)$, and any two cubic rhombi with the same $(n,c)$ differ by an orthogonal transformation of $\R^n$.
If $c<- \frac{1}{5}$ or $c>1$ there is no system of $n$ unit vectors with pairwise inner product $c$.
\end{theorem}
\begin{proof}
The determinant and the range of positivity are Lemma~\ref{lem:spec}.
If $G_n(c)$ is positive definite it has a Cholesky factorization $G_n(c)=V^\T V$ with $V$ invertible, and the columns of $V$ are unit vectors with the required inner products. This gives existence.
If $W$ is another matrix with $W^\T W=G_n(c)$, put $Q=WV^{-1}$. Then $Q^\T Q=V^{-\T}G_n(c)V^{-1}=I$, so $Q$ is orthogonal and $W=QV$. This gives uniqueness.
Finally, the Gram matrix of any system of vectors is positive semidefinite, so a system exists only if both eigenvalues in Lemma~\ref{lem:spec} are nonnegative, that is, only if $- \frac{1}{5}\le c\le1$.
\end{proof}
